% Part: propositional-logic
% Chapter: syntax-and-semantics
% Section: formulas

\documentclass[../../../include/open-logic-section]{subfiles}

\begin{document}

\olfileid[zh]{pl}{syn}{fml}

\olsection{命题逻辑的\usetoken{P}{formula}}

命题逻辑的!!^{formula}由\emph{!!{propositional variable}}\iftag{prvFalse}{\iftag{prvTrue}{，}{和}命题常元~$\lfalse$}{}\iftag{prvTrue}{\iftag{prvFalse}{和}{}命题常元~$\ltrue$}{}借助\emph{逻辑联结词}构成。

\begin{enumerate}
\item !!^a{denumerable}命题变元集~$\PVar$，即 $\Obj p_0$、
  $\Obj p_1$、\dots 等!!{propositional variable}。
\tagitem{prvFalse}{用于\ !!{falsity}的命题常元~$\lfalse$。}{}
\tagitem{prvTrue}{用于\ !!{truth}的命题常元~$\ltrue$。}{}
\item 逻辑联结词：
  \startycommalist
  \iftag{prvNot}{\ycomma $\lnot$（否定）}{}%
  \iftag{prvAnd}{\ycomma $\land$（合取）}{}%
  \iftag{prvOr}{\ycomma $\lor$（析取）}{}%
  \iftag{prvIf}{\ycomma $\lif$（!!{conditional}）}{}%
  \iftag{prvIff}{\ycomma $\liff$（!!{biconditional}）}{}%
\item 标点符号：左括号「(」、右括号「)」以及逗号。
\end{enumerate}

我们将上述命题逻辑的语言记作 $\Lang L_0$。

\iftag{defNot,defOr,defAnd,defIf,defIff,defTrue,defFalse,defEx,defAll}{%
除上面引入的初始联结词外，我们还使用下列\emph{定义的}符号：
\startycommalist
  \iftag{defNot}{\ycomma $\lnot$（否定）}{}%
  \iftag{defAnd}{\ycomma $\land$（合取）}{}%
  \iftag{defOr}{\ycomma $\lor$（析取）}{}%
  \iftag{defIf}{\ycomma $\lif$（!!{conditional}）}{}%
  \iftag{defIff}{\ycomma $\liff$（!!{biconditional}）}{}%
  \iftag{defFalse}{\ycomma $\lfalse$（!!{falsity}）}{}%
  \iftag{defTrue}{\ycomma $\ltrue$（!!{truth}）}}{}。

\begin{tagblock}{defNot,defOr,defAnd,defIf,defIff,defTrue,defFalse,defEx,defAll}
\begin{explain}
定义的符号并不是正式属于该语言的符号，而是作为一种非正式缩写引入的：它使我们得以缩写那些若只用初始符号就会变得相当长的公式。这显然是一种好处。然而更大的好处在于，证明会变得更短。如果一个符号是初始的，那么它在证明中就必须单独处理。因此，初始符号越多，我们的证明就越长。
\end{explain}
\end{tagblock}

% Alternate symbols

\begin{intro}
你可能熟悉与我们上面所用不同的术语和符号。逻辑教科书（和老师们）通常使用 $\sim$、$\neg$ 以及~$!$ 来表示「否定」，用 $\wedge$、$\cdot$ 和~$\&$ 来表示「合取」。「条件句」或「蕴含」常用的符号是 $\rightarrow$、$\Rightarrow$ 和~$\supset$。
\iftag{prvIff,defIff}{「双条件」「双蕴含」或「（实质）等值」的符号是 $\leftrightarrow$、
  $\Leftrightarrow$ 和~$\equiv$。}{}
\iftag{prvFalse,defFalse}{$\lfalse$ 符号有不同的叫法：\textit{falsity}、\textit{falsum}、\textit{absurdity} 或 \textit{bottom}。}{}
\iftag{prvTrue,defTrue}{$\ltrue$ 符号有不同的叫法：\textit{truth}、\textit{verum} 或 \textit{top}。}{}
\end{intro}

\begin{defn}[公式]
\ollabel{defn:formulas}
命题逻辑的\emph{!!{formula}}的集合~$\Frm[L_0]$ 归纳定义如下：
\begin{enumerate}
\tagitem{prvFalse}{$\lfalse$ 是一个原子!!{formula}。}{}

\tagitem{prvTrue}{$\ltrue$ 是一个原子!!{formula}。}{}

\item 每个!!{propositional variable}~$\Obj p_i$ 都是一个原子!!{formula}。

\tagitem{prvNot}{若 $!A$ 是一个!!a{formula}，则 $\lnot !A$ 是一个!!a{formula}。}{}

\tagitem{prvAnd}{若 $!A$ 和 $!B$ 都是!!{formula}，则 $(!A \land
  !B)$ 是一个!!a{formula}。}{}

\tagitem{prvOr}{若 $!A$ 和 $!B$ 都是!!{formula}，则 $(!A \lor !B)$
  是一个!!a{formula}。}{}

\tagitem{prvIf}{若 $!A$ 和 $!B$ 都是!!{formula}，则 $(!A \lif !B)$
  是一个!!a{formula}。}{}

\tagitem{prvIff}{若 $!A$ 和 $!B$ 都是!!{formula}，则 $(!A \liff !B)$
  是一个!!a{formula}。}{}

\tagitem{limitClause}{除此之外没有别的!!a{formula}。}{}
\end{enumerate}
\end{defn}

\begin{explain}
!!{formula}的定义是一个
\emph{归纳定义}。本质上，我们是在无穷多个阶段中构造!!{formula}的集合的。在初始阶段，我们判定所有原子公式都是公式；这对应于定义中的最初几个情形，即
\iftag{prvTrue}{$\ltrue$， }{}%
\iftag{prvFalse}{$\lfalse$， }{}%
$\Obj p_i$ 的情形。「原子!!{formula}」
因此意指任何这种形式的!!{formula}。

定义中的其他情形给出了由已经构造出的!!{formula}构造新的!!{formula}的规则。在第二阶段，我们可以用它们由原子!!{formula}构造!!{formula}。在第三阶段，我们由原子公式以及在第二阶段得到的那些公式构造新的公式，如此等等。一个!!{formula}就是在某一这样的阶段最终构造出的任何东西，此外没有别的。
\end{explain}

当我们用二元联结词~$\ast$ 由 $!B$、$!C$
构造公式 $(!B \ast !C)$ 时，我们常常省去最外面的一对括号，只简单地写成~$!B \ast !C$。

\begin{tagblock}{defNot,defOr,defAnd,defIf,defIff,defTrue,defFalse,defEx,defAll}
\begin{defn}
使用定义的算子构造的公式应按如下方式理解：

\begin{tagenumerate}{defTrue,defFalse,defNot,defOr,defAnd,defIf,defIff,defEx,defAll}

\tagitem{defTrue}{$\ltrue$ 缩写为
  \iftag{prvFalse}{$\lnot\lfalse$}{$(!A \lor \lnot !A)$，其中 $!A$ 是某个
    固定的原子!!{formula}}。}{}

\tagitem{defFalse}{$\lfalse$ 缩写为
  \iftag{prvTrue}{$\lnot\ltrue$}{$(!A \land \lnot !A)$，其中 $!A$ 是某个
    固定的原子!!{formula}}。}{}

\tagitem{defNot}{$\lnot !A$ 缩写为 $!A \lif \lfalse$。}{}

\tagitem{defOr}{$!A \lor !B$ 缩写为
  \iftag{prvAnd}{$\lnot(\lnot !A \land \lnot !B)$}{$\lnot !A \lif
    !B$}。}{}

\tagitem{defAnd}{$!A \land !B$ 缩写为
  \iftag{prvOr}{$\lnot(\lnot !A \lor \lnot !B)$}{$\lnot (!A \lif
    \lnot !B)$}。}{}

\tagitem{defIf}{$!A \lif !B$ 缩写为
  \iftag{prvOr}{$\lnot !A \lor !B)$}{$\lnot (!A \land \lnot !B)$}。}{}

\tagitem{defIff}{$!A \liff !B$ 缩写为 $(!A \lif !B) \land (!B
  \lif !A)$。}{}
\end{tagenumerate}
\end{defn}
\end{tagblock}

\begin{defn}[句法等同]
符号 $\ident$ 表示符号串之间的句法等同，即 $!A \ident !B$ 当且仅当 $!A$ 和 $!B$ 是长度相同且在每一位上含有相同符号的符号串。
\end{defn}

$\ident$ 的两侧可以是经连接得到的字符串，例如 $!A \ident (!B \lor !C)$ 的意思是：符号串~$!A$ 与将左括号、字符串 $!B$、$\lor$ 符号、字符串~$!C$ 以及右括号依次连接所得到的字符串是同一个字符串。若如此，
那么我们就知道 $!A$ 的第一个符号是左括号，$!A$ 在第二个符号起含有 $!B$ 作为子字符串，该子字符串之后是~$\lor$，等等。

\end{document}
